% =============================================================================
% Resumen y abstract (un párrafo de 250 palabras como máximo, con 4 palabras clave).
% Versión de avance (MMC102, 7-nov-2026): resume solo lo que sostienen los capítulos 3 a 6.
% Las cifras de datos entran por macros de cifras.tex.
% =============================================================================
\chapter*{Resumen}
\addcontentsline{toc}{chapter}{Resumen}
\markboth{Resumen}{Resumen}

Esta tesis formaliza y aplica una metodología exploratoria evolutiva para detectar, fechar y
caracterizar los regímenes del ciclo financiero peruano con \nSeries{} series mensuales del Banco
Central de Reserva del Perú, de \fechaInicioNiveles{} a \fechaFinDatos{}, analizadas sin
transformar, según el principio de «datos primero». La cadena integra el análisis de componentes
principales (PCA), la Clasificación Factorial Jerárquica (HFC), el análisis evolutivo de
componentes principales (EPCA) y la partición óptima de Fisher; los autoencoders y las redes
LSTM se formalizan como un complemento anidado en el PCA. Como aporte formal, se demuestran
propiedades de la correlación ponderada por ventana, del tamaño efectivo de los pesos, del
coeficiente RV de cargas y de la partición con longitud mínima, y se propone un criterio para el
ancho de ventana del EPCA, basado en el RV entre ventanas adyacentes no solapadas y contrastado
con un nulo de estructura constante. Los resultados de avance muestran un primer eje dominado por
la tendencia común de profundización financiera, que no supera lo esperable en series
independientes con su propia tendencia, y un eje de tasas de interés cuyo comovimiento sí lo
supera. La HFC separa de forma estable tres bloques: crédito, liquidez y actividad; tasas de
interés; y riesgo de mercado. Con las series sin transformar, el criterio de ventana no distingue
los cambios de estructura de lo esperable bajo estructura constante, y el número de regímenes
fechados depende de la regla que lo fija. Las conclusiones son preliminares y quedan sujetas a
la discusión con el asesor.

\medskip
\textbf{Palabras clave:} ciclo financiero; Clasificación Factorial Jerárquica; análisis evolutivo
de componentes principales; detección de regímenes.

\chapter*{Abstract}
\addcontentsline{toc}{chapter}{Abstract}
\markboth{Abstract}{Abstract}

\begin{otherlanguage}{english}
This thesis formalizes and applies an evolutionary exploratory methodology to detect, date and
characterize the regimes of the Peruvian financial cycle, using \nSeries{} monthly series from
the Central Reserve Bank of Peru, from \anioInicioDatos{} to \anioFinDatos{}, analyzed without
transformation, following the ``data first'' principle. The chain combines principal component
analysis (PCA), Hierarchical Factor Classification (HFC), evolutionary principal component
analysis (EPCA) and Fisher's optimal partition; autoencoders and LSTM networks are formalized as
a complement nested in PCA. As a formal contribution, the thesis proves properties of the
window-weighted correlation, of the effective size of the weights, of the RV coefficient of
loadings and of the partition with a minimum segment length, and it proposes a criterion for
the EPCA window width, based on the RV coefficient between adjacent non-overlapping windows and
tested against a constant-structure null. Preliminary results show a first axis dominated by the
common trend of financial deepening, which does not exceed what independent series with their own
trends would produce, and an interest-rate axis whose co-movement does exceed it. HFC stably
separates three blocks: credit, liquidity and activity; interest rates; and market risk. With
untransformed series, the window criterion does not distinguish structural changes from what a
constant structure would produce, and the number of dated regimes depends on the rule that sets
it. The conclusions are preliminary and subject to discussion with the advisor.

\medskip
\textbf{Keywords:} financial cycle; Hierarchical Factor Classification; evolutionary principal
component analysis; regime detection.
\end{otherlanguage}
